\documentclass[11pt,a4paper]{article}
\usepackage[margin=1in]{geometry}
\usepackage{times}
\usepackage{booktabs}
\usepackage{amsmath}
\usepackage{url}
\usepackage[hidelinks]{hyperref}
\usepackage{listings}
\usepackage{xcolor}
\usepackage{enumitem}

\title{Verdict Attestations: A Cross-Registry, Offline-Verifiable\\Ledger for AI Agent Component Supply Chains}
\author{Canari Contributors\\ \small \texttt{github.com/lanbass869-cell/canari}}
\date{November 2026}

\begin{document}
\maketitle

\begin{abstract}
AI agent supply chains face a gap between two existing trust layers: registry-native
identity attestations (Sigstore provenance for npm, PEP~740 for PyPI) prove
\emph{which} code was built, while security scanners enumerate \emph{what} patterns
the code contains. Neither layer produces a portable, offline-verifiable,
time-bound \emph{verdict} about whether a component is safe to install, and no
existing system unifies verdicts across npm, PyPI, the Model Context Protocol (MCP)
registry, and community plugin registries.

We present Canari, a cross-registry verification ledger that produces signed
DSSE/Ed25519 attestation objects binding verdicts to artifact hashes, with
bi-temporal timestamps, supersession chains, and same-author revocation propagation.
Canari does not compete with identity attestations (it consumes them as input) or
with deep scanners (it orchestrates them as adapters); it fills the \emph{verdict
layer} between them. We demonstrate the system with 294 attestations spanning
four registries, a live-verified npm provenance ingestion path
(sigstore-python~4.5.1 against \texttt{express@5.3.0}), and a PyPI native
attestation pipeline that is ready for adoption (measured coverage: 0/40 recent
uploads). We publish all code, data, and schemas as open artifacts.
\end{abstract}

\section{Introduction}

LLM-based agents increasingly rely on third-party components: MCP servers,
agent skills, and plugins from community registries. Two threats dominate this
supply chain. First, the \emph{tool update gap}: trust established for a tool
does not remain valid after implementation updates. Kan et al.~\cite{etude}
demonstrated a 97.6\% attack success rate by embedding malicious functionality
into updates of initially benign tools, bypassing prompt-injection detectors
and static skill scanning entirely. Second, \emph{silent drift}: half of
multi-version MCP servers (51.1\%) changed what they advertise between versions,
40.6\% did so silently, and 4.2\% redirected their remote endpoint to a
different host while keeping the same registry identity~\cite{samename}.

The ecosystem has developed two families of defenses, but a gap remains between
them:

\begin{enumerate}[leftmargin=1.5em]
\item \textbf{Identity attestations} (Sigstore provenance on npm, PEP~740 on PyPI)
bind a release to its builder identity. These prove \emph{which} code runs but not
\emph{what} it does---Kan et al.~\cite{etude} states this explicitly: ``hashes and attestations
certify which code runs, never what it does.'' As of November 2026, 60\% of top-30
npm packages carry Sigstore provenance, but PyPI coverage is effectively zero
(0/30 top packages, 0/40 fresh uploads measured).

\item \textbf{Security scanners} (NVIDIA SkillSpector~\cite{skillspector},
Cisco skill-scanner~\cite{cisco}, Snyk agent-scan~\cite{snyk},
repo-forensics~\cite{repoforensics}) enumerate vulnerability patterns in source
code. The deepest open-source scanner we found ships 1{,}138 regex rules. These
tools produce findings but not signed, portable, time-bound verdicts, and they
do not track drift across versions.
\end{enumerate}

\textbf{The gap.} No existing system produces \emph{signed, offline-verifiable
verdict attestations} that are (a) portable across registries, (b) linked into
supersession chains so history is never lost, (c) propagated through
same-author revocation, and (d) consumable by a client-side install gate.
Closest systems fall short: mcpcheckup~\cite{mcpcheckup} provides signed
monitoring for remote MCP servers only; SIGIL~\cite{sigil} gates skills through
an on-chain registry (imposing chain infrastructure); agtmls~\cite{agtmls}
maintains a signed skill registry with revocation but for a single registry;
and the manifest signing of Jamshidi et al.~\cite{manifest} operates within
MCP pipelines without cross-registry scope.

\section{System Design}

\subsection{Attestation Format}

Each attestation is a JSON payload signed with DSSE~\cite{dsse} over an
Ed25519 keypair. The schema (v0.1) is a domain-specific in-toto predicate.
Core fields:

\begin{itemize}[leftmargin=1.5em]
\item \texttt{subject}: \{name, version, registry, artifact\_sha256\} --- binds the
verdict to exact bytes, not to a mutable name.
\item \texttt{scenario\_sha256}: SHA-256 of the check-list definition; changing
the scanner invalidates prior verdicts.
\item \texttt{checks[]}: per-check verdict (\texttt{pass}/\texttt{fail}/\texttt{skip}/\texttt{unsupported}/\texttt{error}) with tool identifier.
\item \texttt{verdict}: run-level summary (\texttt{pass},
\texttt{pass-with-findings}, \texttt{fail}, \texttt{unsupported}, \texttt{error}).
\item \texttt{supersedes}: digest of the superseded predecessor (supersession chain).
\item \texttt{related\_revocations[]}: same-author / name-variant / dependency
links to revocation events (propagation).
\item \texttt{native\_provenance}: ingested registry-native attestation result
(presence + verification outcome), recorded as an \emph{input signal}, never
as a safety verdict.
\item \texttt{evidence\_completeness}: explicit disclosure of missing evidence
(anti-inflation discipline learned from dsh-testkit: \texttt{unsupported}
$\neq$ \texttt{failed} $\neq$ \texttt{pass}).
\end{itemize}

\subsection{Supersession Chains}

When a component is re-scanned, the new attestation carries \texttt{supersedes}
pointing to the digest of the previous attestation for the same subject. Chains
are append-only: old verdicts are never deleted, so an auditor can reconstruct
the full history of a component's trust state. This mirrors the supersession
semantics of MemTether~\cite{memtether} (a cross-client memory hub) adapted to
supply-chain objects.

\subsection{Same-Author Revocation Propagation}

When a component is revoked (e.g., a registry yank, a scanner \texttt{fail},
or a maintainer removal), Canari queries the state store for all components
sharing the same author, name variants, or declared dependencies, and appends
\texttt{related\_revocations} entries to their next attestations. This is
\emph{annotation}, not judgment: the consumer decides whether a related
revocation affects their install decision.

\subsection{Native Attestation Ingestion}

Canari distinguishes between \emph{identity attestations} (registry-native) and
\emph{verdict attestations} (Canari's own). Native provenance is consumed as an
input signal:

\begin{itemize}[leftmargin=1.5em]
\item \textbf{npm}: the Sigstore bundle (v1, certificate-based) is verified
offline via \texttt{sigstore-python}~4.5.1's \texttt{Verifier.verify\_dsse()}.
Verified live against \texttt{express@5.3.0}: DSSE verification passes,
identity is extracted from the certificate SAN
(\texttt{https://github.com/expressjs/express/.github/workflows/npm-publish.yml@refs/tags/v5.3.0}).
v0.1 bundles (publicKey hint, no certificate chain) are recorded as
\emph{presence-only}; no fake verification is attempted.
\item \textbf{PyPI}: PEP~740 objects are converted to Sigstore Bundles via
\texttt{pypi-attestations}~0.0.31 and verified identically. Measured adoption
is currently zero (0/30 top, 0/40 fresh), but the pipeline is ready.
\end{itemize}

\subsection{Scanner Orchestration}

Canari does not build its own vulnerability scanner. L1--L2 checks (registry
present, version resolved, artifact hash, maintainer set, manifest fields) are
meta-only. L3+ behavioral checks are delegated to external scanners (sentinel,
repo-forensics) via a scanner-adapter interface that records:
scanner name + version, rule count, findings JSON, and a digest of the scanner
source tree (scenario binding). This keeps Canari's core small ($\sim$300 lines
of Python) while allowing deep scans from the best available tools.

\section{Implementation}

Canari is implemented in Python 3.13 with the following dependencies:
\texttt{in-toto} 3.1.0 (DSSE envelopes), \texttt{securesystemslib} 1.6
(Ed25519), \texttt{jsonschema} (schema validation), \texttt{sigstore-python} 4.5.1
(native npm verification), \texttt{pypi-attestations} 0.0.31 (PyPI native).
The codebase is $\sim$1{,}200 lines of Python across 8 modules, with 53 tests
(52 passed, 1 skipped for no live PyPI attestation available).

Adapters exist for four registries:
\begin{itemize}[leftmargin=1.5em]
\item npm (packument + tarball download for hash, search for discovery, native
provenance via \texttt{dist.attestations})
\item PyPI (JSON API + Simple API for full enumeration)
\item MCP Registry (\texttt{registry.modelcontextprotocol.io}, cursor-paginated)
\item dsh.so (market index with 10{,}195 plugins, per-plugin install-test and
security-scan endpoints, CC-BY-4.0 licensed)
\end{itemize}

\section{Evaluation}

\subsection{Coverage}

A multi-registry run produced 294 attestations (npm: 100, PyPI: 94, MCP Registry:
100) with 99.6\% schema integrity and zero hidden errors. Failure counts are
publicly recorded in the feed (dead-letter audit).

\subsection{Native Provenance}

Table~\ref{tab:native} summarizes measured native attestation coverage.

\begin{table}[h]
\centering
\caption{Native attestation coverage (measured 2026-11-08).}
\label{tab:native}
\begin{tabular}{lrr}
\toprule
Registry & Sample & Coverage \\
\midrule
npm (top-30) & 30 & 18/30 (60\%) \\
PyPI (top-30) & 30 & 0/30 (0\%) \\
PyPI (fresh uploads) & 40 & 0/40 (0\%) \\
\bottomrule
\end{tabular}
\end{table}

The npm v1 bundle for \texttt{express@5.3.0} verifies: DSSE signature valid,
identity = \texttt{npm-publish.yml} workflow at tag \texttt{v5.3.0},
tlog entry present. The v0.1 bundle (publicKey hint) is recorded as
presence-only. On PyPI, zero coverage means Canari's PEP 740 pipeline has no
live test target yet; the blank itself is a finding.

\subsection{Drift Cases}

We use the public MCP Registry Drift Panel~\cite{drift89} (Zenodo 21798111,
CC-BY-4.0) as baseline data. The panel records 18{,}966 distinct servers over
88.6 days. Key statistics: 75.2\% of servers never change; the most active 5\%
generate 61\% of all change events; drift-ranked re-auditing at a 5\% budget
catches only $\sim$20\% of drifts. Canari's contribution is the \emph{verdict
layer} on top of this measurement data: for each drift event, Canari produces a
signed re-attestation whose \texttt{supersedes} pointer makes the drift
auditable and reversible.

\subsection{Propagation}

Five test cases verify the same-author revocation propagation logic:
same-author detection, name-variant matching, no self-revocation, empty set,
and relation-field correctness. The propagation is annotation-only; the
consumer's install gate decides.

\subsection{Honest Scoring}

Canari publishes a four-dimension score with every run:
coverage (attested / enumerated), freshness (within TTL), integrity (signature
+ schema pass rate), and failure transparency (error counts). All scores
include sample size, calculation method, and known distortions --- a score
without these is a lie.

\section{Related Work}

Table~\ref{tab:related} compares Canari with the closest systems.

\begin{table*}[t]
\centering
\caption{Canari vs.\ the closest systems (verified against source code, not marketing).}
\label{tab:related}
\small
\begin{tabular}{lcccccc}
\toprule
System & Cross-registry & Signed & Offline & Supersession & Propagation & Gate \\
\midrule
mcpcheckup~\cite{mcpcheckup} & no (MCP only) & yes & yes & no & no & no \\
SIGIL~\cite{sigil} & no (skills) & yes (chain) & no (chain) & no & no & yes (SVL) \\
agtmls~\cite{agtmls} & no (own registry) & yes & yes & no & no (advisory) & lockfile \\
Manifest signing~\cite{manifest} & no (MCP pipeline) & yes & yes & no & no & pipeline \\
Native provenance (npm/PyPI) & per-registry & yes & yes & no & no & advisory \\
SkillSpector~\cite{skillspector} & no & no & -- & no & no & no \\
repo-forensics~\cite{repoforensics} & no & no & -- & no & no & no \\
trustcheck~\cite{trustcheck} & no (PyPI) & yes (PEP 740) & yes & no & no & no \\
\midrule
Canari & \textbf{yes} & \textbf{yes} & \textbf{yes} & \textbf{yes} & \textbf{yes} & \textbf{yes} \\
\bottomrule
\end{tabular}
\end{table*}

\section{Limitations}

\begin{enumerate}[leftmargin=1.5em]
\item \textbf{Scan depth zero.} Canari's built-in checks are meta-only (L1--L2).
Deep vulnerability detection requires external scanners (L3+). This is by design:
Canari is a ledger, not a scanner. The adapter interface makes this composable.
\item \textbf{Single signer.} All attestations in the current deployment are
signed by one key. A production system would need multi-signer support or a
transparency log for the signer keys themselves.
\item \textbf{PyPI blank.} Zero native attestation coverage means the PyPI
ingestion path has no live verification target. The pipeline is ready but
untested against real data.
\item \textbf{MCP Registry count is a floor.} The live API caps pagination at
400 pages ($\sim$40{,}000 entries including versions); distinct-server count
requires resuming pagination with backoff.
\item \textbf{npm v0.1 bundles unverifiable with off-the-shelf tools.}
sigstore-python~4.5.1's \texttt{Bundle} model requires a certificate chain;
v0.1 publicKey-hint bundles are recorded as presence-only. A custom DSSE
verifier could close this gap.
\end{enumerate}

\section{Conclusion}

Canari fills a specific gap in the AI agent supply chain: the verdict layer
between identity attestations (which prove who built the code) and deep
scanners (which enumerate what patterns the code contains). Its contribution
is the combination of cross-registry scope, offline-verifiable signed
attestations, supersession chains, same-author revocation propagation, and a
consumer-side install gate --- four capabilities that no existing system
provides together. All code, schemas, and data are open artifacts.

\bibliographystyle{plainnat}
\begin{thebibliography}{99}

\bibitem{etude} S.~Kan, Y.~Qi, Y.~Han, X.~Luo, G.~Xiong, X.~Gao.
Beyond Prompt Injection: The Tool Update Gap in LLM Agent Security.
\emph{arXiv:2611.01976}, 2026.

\bibitem{samename} O.~Kraishan.
Same Name, Different Server: A Security Census of Silent Drift in the Model Context Protocol Ecosystem.
\emph{arXiv:2609.14119}, 2026.

\bibitem{drift89} G.~Bharti.
Registry Descriptions Go Stale Unevenly: An 89-Day Measurement of MCP Drift.
\emph{arXiv:2608.00997}, 2026. Data: \texttt{zenodo.org/records/21798111}.

\bibitem{mcpcheckup} MCP Checkup.
\texttt{github.com/mcpcheckup/checks}, 2026. Live: \texttt{mcpcheckup.com}.

\bibitem{sigil} T.~Shen et al.
Sealing the Audit-Runtime Gap for LLM Skills (SIGIL).
\emph{arXiv:2605.05274}, 2026.

\bibitem{agtmls} S.~Rousseau.
AgtMLS: Signed, Audited Agent Skills Registry.
\texttt{github.com/sebastienrousseau/agtmls}, 2026.

\bibitem{manifest} S.~Jamshidi et al.
Verifiable Manifest Signing and Transparency Enforcement for Secure MCP-Based LLM Pipelines.
\emph{arXiv:2601.23132}, 2026.

\bibitem{skillspector} NVIDIA.
SkillSpector: Security Scanner for AI Agent Skills.
\texttt{github.com/NVIDIA/SkillSpector}, 2026.

\bibitem{cisco} Cisco AI Defense.
Skill Scanner / MCP Scanner.
\texttt{github.com/cisco-ai-defense}, 2026.

\bibitem{snyk} Snyk.
Agent Scan.
\texttt{github.com/snyk/agent-scan}, 2026.

\bibitem{repoforensics} A.~Greensh.
repo-forensics: Offline Security Scanner for AI-Agent Repos.
\texttt{github.com/alexgreensh/repo-forensics}, 2026.

\bibitem{trustcheck} K.~Duraphe.
trustcheck: Verify PyPI Package Attestations.
\texttt{github.com/kapilduraphe/trustcheck}, 2026.

\bibitem{dsse} M.~J.~F.~Rehmani et al.
DSSE: Dead Simple Signing Envelope.
\texttt{github.com/secure-systems-lab/dsse}, 2020.

\bibitem{memtether} Canari Contributors.
MemTether: A Cross-Client AI Memory Hub.
Internal reference, 2026.

\bibitem{mcpeco} K.~Tounoussi et al.
MCP-Eco: A Linked Dataset of MCP Servers, Consumers, and Development Activity.
\emph{arXiv:2611.06897}, 2026. Data: \texttt{huggingface.co/datasets/mcp-research-alliance/mcp-ecosystem}.

\bibitem{probeid} G.~Chen et al.
ProbeID: Verifying MCP Tool Identity for LLM Agents under Registry Drift.
\emph{arXiv:2610.33633}, 2026.

\end{thebibliography}

\end{document}
